% allgem. Dokumentenformat
\documentclass[a4paper,12pt,headsepline]{scrartcl}
%Variablen welche innerhalb der gesamten Arbeit zur Verfügung stehen sollen
\newcommand{\titleDocument}{Exposee zur Bachelorthesis}
\newcommand{\subjectDocument}{im Studiengang Duales Studium Informatik}
\newcommand{\domain}{Commerz Systems GmbH}

% weitere Pakete
% Grafiken aus PNG Dateien einbinden
\usepackage{graphicx}

% Text neben Grafik
%\usepackage{wrapfig}

% Text neben Grafik
%\usepackage{censor}

% Deutsche Sonderzeichen benutzen 
%\usepackage{ngerman}

% deutsche Silbentrennung
\usepackage[ngerman]{babel}

% Eurozeichen einbinden
\usepackage[right]{eurosym}

% Umlaute unter UTF8 nutzen
\usepackage[utf8]{inputenc}

% Zeichenencoding
\usepackage[T1]{fontenc}

\usepackage{lmodern}
\usepackage{fix-cm}

% floatende Bilder ermöglichen
%\usepackage{floatflt}

% mehrseitige Tabellen ermöglichen
\usepackage{longtable}

% Unterstützung für Schriftarten
%\newcommand{\changefont}[3]{ 
%\fontfamily{#1} \fontseries{#2} \fontshape{#3} \selectfont}

% Packet für Seitenrandabständex und Einstellung für Seitenränder
\usepackage{geometry}
\geometry{left=3.5cm, right=2cm, top=2.5cm, bottom=2cm}

% Paket für Boxen im Text
%\usepackage{fancybox}

% bricht lange URLs "schoen" um
\usepackage[hyphens,obeyspaces,spaces]{url}

% Paket für Textfarben
\usepackage{color}

% Mathematische Symbole importieren
\usepackage{amssymb}

% auf jeder Seite eine Überschrift (alt, zentriert)
%\pagestyle{headings}

% erzeugt Inhaltsverzeichnis mit Querverweisen zu den Kapiteln (PDF Version)
\usepackage[bookmarksnumbered,pdftitle={\titleDocument},hyperfootnotes=false]{hyperref} 
%\hypersetup{colorlinks, citecolor=red, linkcolor=blue, urlcolor=black}
%\hypersetup{colorlinks, citecolor=black, linkcolor= black, urlcolor=black}

% neue Kopfzeilen mit fancypaket
%\usepackage{fancyhdr} %Paket laden
%\pagestyle{fancy} %eigener Seitenstil
%\fancyhf{} %alle Kopf- und Fußzeilenfelder bereinigen
%\fancyhead[L]{\nouppercase{\leftmark}} %Kopfzeile links
%\fancyhead[C]{} %zentrierte Kopfzeile
%\fancyhead[R]{\thepage} %Kopfzeile rechts
%\renewcommand{\headrulewidth}{0.4pt} %obere Trennlinie
%\fancyfoot[C]{\thepage} %Seitennummer
%\renewcommand{\footrulewidth}{0.4pt} %untere Trennlinie

% für Tabellen
\usepackage{array}

% Runde Klammern für Zitate
%\usepackage[numbers,round]{natbib}

% Festlegung Art der Zitierung - Havardmethode: Abkuerzung Autor + Jahr
\bibliographystyle{alphadin}

% Schaltet den zusätzlichen Zwischenraum ab, den LaTeX normalerweise nach einem Satzzeichen einfügt.
\frenchspacing

% Paket für Zeilenabstand
\usepackage{setspace}

% für Bildbezeichner
%\usepackage{capt-of}

% für Stichwortverzeichnis
\usepackage{makeidx}

% für Listings
% \usepackage{listings}
% \lstset{numbers=left, numberstyle=\tiny, numbersep=5pt, keywordstyle=\color{black}\bfseries, stringstyle=\ttfamily,showstringspaces=false,basicstyle=\footnotesize,captionpos=b}
% \lstset{language=java}

\usepackage{listings}
\usepackage{color}

% Blocksatz?
%\usepackage[none]{hyphenat}

\definecolor{mygreen}{rgb}{0,0.6,0}
\definecolor{mygray}{rgb}{0.5,0.5,0.5}
\definecolor{mymauve}{rgb}{0.58,0,0.82}


% Indexerstellung
\makeindex

% Abkürzungsverzeichnis
%\usepackage[german]{nomencl}
%\let\abbrev\nomenclature

% Abkürzungsverzeichnis LiveTex Version
%\renewcommand{\nomname}{Abkürzungsverzeichnis}
%\setlength{\nomlabelwidth}{0.01\hsize}
%\renewcommand{\nomlabel}[1]{#1 \dotfill}
%\setlength{\nomitemsep}{-\parsep}
%\makenomenclature
%\makeglossary

% Abkürzungsverzeichnis TeTEX Version
%\usepackage[german]{nomencl}
%\makenomenclature
%%\makeglossary
%\renewcommand{\nomname}{Abkürzungsverzeichnis}
%\setlength{\nomlabelwidth}{.25\hsize}
%\renewcommand{\nomlabel}[1]{#1 \dotfill}
%\setlength{\nomitemsep}{-\parsep}

% Disable single lines at the start of a paragraph (Schusterjungen)
\clubpenalty = 10000
% Disable single lines at the end of a paragraph (Hurenkinder)
\widowpenalty = 10000
\displaywidowpenalty = 10000

\begin{document}
% hier werden die Trennvorschläge inkludiert
%hier müssen alle Wörter rein, welche Latex von sich auch nicht korrekt trennt bzw. bei denen man die genaue Trennung vorgeben möchte
\hyphenation{
Film-pro-du-zen-ten
Lux-em-burg
Soft-ware-bau-steins
zeit-in-ten-siv
}

% Blocksatz Einstellung
\hyphenpenalty=5000
\exhyphenpenalty=5000
\sloppy

%Schriftart Helvetica
%\changefont{phv}{m}{n}

% Leere Seite am Anfang

\thispagestyle{empty} % erzeugt Seite ohne Kopf- / Fusszeile
\section*{ }

% Titelseite %
\thispagestyle{empty}

%\begin{figure}[t]
% \includegraphics[width=0.6\textwidth]{abb/fh_koeln_logo}
%\end{figure}

\begin{figure}[t]
 \centering
 \includegraphics[width=0.5\textwidth]{abb/logo}
\end{figure}

%\begin{tabular}{lr}
%	\includegraphics[width=0.5\textwidth]{abb/logo}	&	\includegraphics[width=0.5\textwidth]{abb/logo_cs}
%\end{tabular}

\begin{center}
\Large{Hochschule Bremen}\\
\Large{Fakultät für Elektrotechnik und Informatik}\\
\end{center}



\begin{verbatim}


\end{verbatim}
\begin{center}
\onehalfspacing
\textbf{\LARGE{Untersuchung und Evaluierung der \\ Möglichkeiten für die Realisierung \\ automatisierter Tests für \\AngularJS-Webanwendungen \\}}
\singlespacing
\begin{verbatim}



\end{verbatim}
\textbf{\titleDocument \\ zur Erlangung des akademischen Grades \\ Bachelor of Science (B.\,Sc.) \\ im Dualen Studiengang Informatik}
\end{center}
\begin{verbatim}






\end{verbatim}
\begin{center}
\begin{tabular}{llll}
\textbf{Autor} & & Nikolas Schreck <nikolasschreck@gmail.com> & \\
& & \\
\textbf{Version vom:} & & \today &\\
& & \\
\textbf{Erstprüfer} & & Prof. Dr.-Ing. Heide-Rose Vatterrott&\\
\textbf{Zweitprüfer} & & Dipl.-Inf. Jochen Schwitzing&\\
\end{tabular}
\end{center}

% römische Numerierung
%\pagenumbering{arabic}

% 1.5 facher Zeilenabstand
\onehalfspacing

% Sperrvermerk
%\section*{Sperrvermerk}
\textcolor{red}{
Die vorliegende Prüfungsarbeit enthält vertrauliche Daten der \domain, die der Geheimhaltung 
unterliegen. Die Prüfungsarbeit wird an der Hochschule Bremen ausschließlich solchen 
Personen zugänglich gemacht, die mit der Abwicklung des Prüfungsverfahrens betraut sind und 
zur Verschwiegenheit verpflichtet sind. Es wird darauf hingewiesen, dass, sofern der 
Verfasser die Bewertung seiner Arbeit angreift, die Arbeit gegebenenfalls dem 
Widerspruchsausschuss zugeleitet werden muss, wobei die Mitglieder des Widerspruchsausschusses zur Verschwiegenheit verpflichtet sind. Wird die Bewertung der Arbeit 
gerichtlich angegriffen, so ist die Arbeit als Teil des Verwaltungsvorgangs dem Gericht zu 
übermitteln. Veröffentlichung und Vervielfältigung der vorliegenden Prüfungsarbeit – auch 
nur auszugsweise und gleich in welcher Form – bedürfen der schriftlichen Genehmigung der \domain.
}

% Einleitung / Abstract
%\section*{Zusammenfassung}
In dieser \titleDocument{} werden die Möglichkeiten zur Realisierung automatisierter Tests von AngularJS-Anwendungen untersucht. Hierbei werden speziell die Erfordernisse des Commerzbank-Konzern beachtet.

Zunächst wird eine Einführung in die allgemeinen Themen Testen, Angular und Node.js gegeben und Anforderungen an Testframework und -tools definiert. Anschließend werden die verschiedene Programme untersucht und mit Blick auf die Erfüllung der aufgestellten Anforderungen bewertet. Eine Kombination von Karma, Mocha, Chai und Protractor, sowie Sinon, ngMock und Istanbul, wird für den Einsatz im Commerzbank-Konzern als am geeignetsten betrachtet und ausgewählt.

Eine prototypische Realisierung von Tests in einem exemplarischen, realen Projekt mit der Auswahl zeigt, dass die Realisierung von Komponententests mit Karma, Mocha und Chai problemlos funktioniert und einfach anzuwenden ist. Es zeigt sich jedoch auch, dass eine Durchführung von automatisierten End-To-End-Tests an Commerzbank-Arbeitsplätzen aufgrund verschiedener regulatorischer Vorgaben und Einschränkungen nicht möglich ist und die ausgewählte Software Protractor für diesen Zweck nicht funktioniert.

Abschließend wird ein Ausblick auf mögliche Verbesserungen bei automatisierten Tests durch den Einsatz neuerer Versionen von Angular vorgenommen und ein positives Fazit gezogen.

% einfacher Zeilenabstand
%\singlespacing

% Inhaltsverzeichnis anzeigen
%\newpage
%\tableofcontents

% das Abbildungsverzeichnis
% Abbildungsverzeichnis soll im Inhaltsverzeichnis auftauchen
%\addcontentsline{toc}{section}{Abbildungsverzeichnis}
% Abbildungsverzeichnis endgueltig anzeigen
%\listoffigures

% das Tabellenverzeichnis
% Abbildungsverzeichnis soll im Inhaltsverzeichnis auftauchen
%\addcontentsline{toc}{section}{Tabellenverzeichnis}
% \fancyhead[L]{Abbildungsverzeichnis / Abkürzungsverzeichnis} %Kopfzeile links
% Abbildungsverzeichnis endgueltig anzeigen
%\listoftables

%% WORKAROUND für Listings
%\makeatletter% --> De-TeX-FAQ
%\renewcommand*{\lstlistoflistings}{%
%  \begingroup
%    \if@twocolumn
%      \@restonecoltrue\onecolumn
%    \else
%      \@restonecolfalse
%    \fi
%    \lol@heading
%    \setlength{\parskip}{\z@}%
%    \setlength{\parindent}{\z@}%
%    \setlength{\parfillskip}{\z@ \@plus 1fil}%
%    \@starttoc{lol}%
%    \if@restonecol\twocolumn\fi
%  \endgroup
%}
%\makeatother% --> \makeatletter
% das Listingverzeichnis
%\newpage
% Listingverzeichnis soll im Inhaltsverzeichnis auftauchen
%\addcontentsline{toc}{section}{Listingverzeichnis}
%\fancyhead[L]{Abbildungs- / Tabellen- / Listingverzeichnis} %Kopfzeile links
%\renewcommand{\lstlistlistingname}{Listingverzeichnis}
%\lstlistoflistings
%%%%

% das Abkürzungsverzeichnis
%\newpage
% Abkürzungsverzeichnis soll im Inhaltsverzeichnis auftauchen
%\addcontentsline{toc}{section}{Abkürzungsverzeichnis}
% das Abkürzungsverzeichnis entgültige Ausgeben
%\fancyhead[L]{Abkürzungsverzeichnis} %Kopfzeile links
%\nomenclature{API}{Application Programming Interface}
\nomenclature{JSON}{JavaScript Object Notation}
\nomenclature{BaFin}{Bundesanstalt für Finanzdienstleistungsaufsicht}
\nomenclature{jQuery}{jQuery ist eine JavaScript-Bibliothek für DOM-Manipulation, Event-Handling und Animation.}
\nomenclature{Whitebox-Test}{Ein Whitebox-Test ist ein Test, welcher unter Kenntniss der inneren Funktionsweise entwickelt wird.}
\nomenclature{REST}{Representational State Transfer}
\nomenclature{Mock}{Ein Platzhalter-Objekt, welches als Attrappe verwendet wird.}
\nomenclature{MVVM}{Model-View-ViewModel, ein Entwurfsmuster}
\nomenclature{MVC}{Model-View-Controller, ein Entwurfsmuster}
\nomenclature{HTTP}{Hypertext Transfer Protocol}
\nomenclature{WebSocket}{WebSocket ist eine bidirektionale Erweiterung von HTTP.}
\nomenclature{DOM}{Document Object Model}
\nomenclature{falsy}{Ein in JavaScript als false angesehener Wert: false, 0, \glqq\grqq (leerer String), null, undefined und NaN}
\nomenclature{truthy}{Ein in JavaScript als true angesehener Wert: alle nicht falsy Werte, auch \glqq 0\grqq, \glqq false\grqq{} oder leere Funktionen, Arrays, Objekte}

%\printnomenclature

% Definiert Stegbreite bei zweispaltigem Layout
\setlength{\columnsep}{25pt}

%%%%%%% EINLEITUNG %%%%%%%%%%%%
%\twocolumn
%\fancyhead[L]{\nouppercase{\leftmark}} %Kopfzeile links

% 1,5 facher Zeilenabstand
\onehalfspacing

% einzelne Kapitel
\section{Motivation}\label{einleitung}
Die Commerzbank AG ist Deutschlands zweitgrößtes Finanzinstitut\cite[][2]{handelsblatt:commerzbank}. Als solches stellen sich für ihre IT besondere Herausforderungen an Daten-, Ausfallsicherheit und Stabilität. Regulatorische Vorgaben durch die BaFin\cite{bafin-banken}, interne Programmierrichtlinien\cite{coba-programmierrichtlinienAllgemein, coba-programmierrichtlinienJavaScript}, das Book of Standards \cite{coba-bookOfStandards} und Weitere führen zu einem trägen und wenig innovativem IT-Umfeld. Somit ist es nicht überraschend, dass zur Entwicklung von Webanwendungen noch immer auf alte und etablierte Technologien wie jQuery und JavaServer Pages, welche von Oracle als veraltet erklärt wurden \cite{jsp-deprecated}, zurückgegriffen wird\todo{Quelle!}.

Viele Unternehmen nutzen bereits seit einigen Jahren Angular als Technologiebasis für den clientseitigen Teil von Webanwendungen, beispielsweise \textit{Gmail}, \textit{PayPal} oder \textit{Youtube}\cite{angular-sites}. In der Commerzbank wurde Angular bisher nicht berücksichtigt; erst in jüngerer Vergangenheit wird es in vereinzelten Projekten eingesetzt. Hierbei werden automatisierte Entwicklertests in JavaScript aufgrund von Unwissenheit über die Möglichkeiten meist stiefmütterlich behandelt. Die Anwendungen werden stattdessen per Hand im Webbrowser getestet. Die sich hieraus ergebende Testabdeckung steht im Gegensatz zu den Anforderungen an Sicherheit und Stabilität im Bankenumfeld.

In dieser \titleDocument~ sollen daher die Möglichkeiten zur Realisierung von automatisierten Tests in AngularJS-Webanwendungen untersucht werden.

%\newpage
%\section{Grundlagen}\label{grundlagen}
\subsection{Test}
Der Test von Software dient dazu, mögliche Fehler aufzudecken und dadurch die Qualität zu erhöhen. Der Nachweis von Fehlerfreiheit ist unmöglich, daher muss der Testaufwand verhältnismäßig zum Ergebnis sein\cite[][14\psqq]{spillner}. 

Beim Testen werden üblicherweise vier Teststufen unterschieden\cite[][42\psq]{spillner}:
\begin{itemize}
	\item Komponententest
	\item Integrationstest
	\item Systemtest
	\item Abnahmetest
\end{itemize}
Bei Komponenten-, Integrations- sowie bedingt bei Systemstests handelt es sich um Entwicklertests, weshalb diese im Rahmen dieser \titleDocument{} relevant sind. Der Abnahmetest wird daher nicht betrachtet.

\subsubsection{Komponententest}
Ein Komponententest überprüft die einzelnen Bausteine der entwickelten Software erstmalig und unabhängig von anderen Bausteinen. Es wird überprüft, ob die Komponente den Anforderungen sowie dem definierten Softwaredesign entspricht. Außerdem kann auch der Quellcode analysiert und zur Erstellung der Testfälle herangezogen werden; dann handelt es sich um einen \textit{Whitebox-Test}.\cite[][44]{spillner} In JavaScript ist die kleinste testbare Komponente üblicherweise eine Funktion\cite{smashing-unit}. Die zu testende Komponente muss nicht zwingend atomar sein, d.\,h. die Funktion kann aus weiteren Funktionen zusammengesetzt sein, jedoch sollte nur die komponenteninterne Funktionsweise getestet werden\cite[][45]{spillner}. Beim Test von AngularJS-Anwendungen sind die Komponenten beispielsweise Controller, Services, Direktiven und Filter.

Ein Komponententest hat spezifische Testziele. Das Wichtigste ist die Sicherstellung, dass die Funktion der Anforderung in der Spezifikation entspricht. Hierdurch wird die Komponente wie gefordert mit anderen Komponenten zusammenarbeiten und somit in die Gesamtsoftware integriert werden können. Wichtig ist auch der Test auf Robustheit: Bei falschem Aufruf, also einem Verstoß gegen die Vorbedingungen, sollte die Komponente sinnvoll reagieren und den Fehler abfangen. Testfälle lassen sich in Positiv- und Negativtests unterteilen: Positivtests sind die Überprüfung von vorgesehenem Verhalten der Komponente, Negativtests der Test von nicht vorgesehenen, unzulässigen oder explizit ausgeschlossenen Sonderfällen.\cite[][48]{spillner}.

Im Komponententest können auch nicht funktionale Qualitätseigenschaften getestet werden. Zu nennen sind hier beispielsweise Speicherverbrauch oder Antwortzeit, sowie statische Tests \todo{evtl hier Ref auf unterkapitel}auf Wartbarkeit, wie beispielsweise vorhandene Quelltextkommentare oder die Einhaltung von Programmierrichtlinien.\cite[][49\psq]{spillner}.

\subsubsection{Integrationstest}
Der Integrationstest folgt nach den Komponententests und basiert auf getesteten Komponenten. Diese Komponenten werden zu größeren Komponenten oder Teilsystemen zusammengesetzt. Der Integrationstest dient dann dazu zu überprüfen ob alle Einzelteile korrekt zusammenarbeiten und soll Fehler in Schnittstellen und im Zusammenspiel aufdecken. \cite[][52\psq]{spillner} Beispielsweise kann in einem Integrationstest auch die Anbindung an externe Komponenten, wie Datenbanken oder REST-APIs überprüft werden. Diese werden im Komponententest durch \textit{Mocks} ersetzt und emuliert.

Somit ist das aufdecken von Schnittstellenfehlern ein Testziel des Integrationstest, zum Beispiel wegen von der Spezifikation abweichender Schnittstellen. Außerdem ist ein Testziel unerwünschte Wechselwirkungen zwischen den Einzelkomponenten aufzudecken, welche das Zusammenspiel unmöglich machen. \cite[][56]{spillner}

Der Integrationstest ist jedoch kein Ersatz für den Komponententest, da er mit Nachteilen verbunden ist. Es ist schwer bis unmöglich, die tatsächliche Fehlerursache herauszufinden, da oft nicht klar ist in welcher Teilkomponente der Fehler aufgetreten ist, sondern dieser sich nur in einem abweichenden Gesamtverhalten äußert. Manche Fehler werden möglicherweise gar nicht gefunden, da regelmäßig kein vollumfänglicher Zugriff auf Einzelkomponenten besteht.\cite[][57]{spillner}

Es existieren verschiedene Integrationsstrategien, die Auswirkungen auf die Integrationstests haben:\cite[][59\psq]{spillner}
\begin{itemize}
	\item Bei der Top-Down-Integration beginnt der Test mit der obersten Systemkomponente, von der alle Anderen aufgerufen werden. Sukzessive werden die weiteren Komponenten von oben nach unten integriert und getestet, wobei die untergeordneten Komponenten zunächst durch Platzhalter ersetzt werden. Vorteilhaft ist, dass keine aufwändigen Testtreiber zum Aufruf benötigt werden. Jedoch müssen Platzhalterkomponenten implementiert werden, was einen zusätzlichen Overhead beim Test bedeuten kann.
	\item Bei der Bottom-up-Integration werden zunächst die unteren, atomaren Komponenten integriert und getestet. Erst nach und nach werden größere Teilsysteme aus getesteten Komponenten integriert. Der Vorteil hierbei ist, dass keine Platzhalter implementiert werden müssen. Jedoch müssen hier aufwändige Testtreiber erstellt werden, welche die obergeordneten, aufrufenden Komponenten emulieren.
	\item Bei der Ad-hoc-Integration werden Komponenten integriert, sobald sie fertiggestellt sind. Nachteilig hierbei ist, dass sowohl Platzhalter als auch Testtreiber implementiert werden müssen. Allerdings bietet sich ein Zeitgewinn, da jede Komponente so früh wie möglich integriert wird.
	\item Bei der wenig empfehlenswerten Big-Bang-Integration werden alle Komponenten auf einmal integriert. Sie bietet ausschließlich Nachteile: Es wird Zeit verschwendet, da bis zur Fertigstellung der letzten Komponente gewartet wird. Auch treten alle Fehlerwirkungen gesammelt auf, so dass es schwierig ist, die Fehler zu finden.
\end{itemize}

\subsubsection{Systemtest}
Der Systemtest ist der finale Entwicklertest nach den abgeschlossenen Integrationstests. Getestet wird das gesamte System, möglichst in einer produktionsnahen Umgebung. Es sind keine Testtreiber oder Platzhalter mehr vorhanden, stattdessen wird überall die finale Hard- und Software genutzt. Das Testziel ist die Validierung, ob alle funktionalen und nicht-funktionalen Anforderungen erfüllt werden.\cite[][60\psqq]{spillner}

Aufgrund der erforderlichen Produktionsnähe sowie der erforderlichen Datenbasis kann der Systemtest nur bedingt als Entwicklertest gesehen werden\cites[236]{roitzsch}[]{oose}. Jedoch kann und sollte ein einfacher Systemtest auch von Entwicklern durchgeführt werden. Es bietet sich hier die Durchführung von End-To-End-Tests an, mittels derer das System von der Benutzeroberfläche bis zur untersten Komponentenschicht getestet werden kann\cite{softwaresanierung}.

\subsection{AngularJS}
AngularJS ist ein von Google Inc. ins Leben gerufenes JavaScript-Framework zur Entwicklung von clientseitigen Webanwendungen\cite{angular-faq}. Der Quellcode von AngularJS steht auf Github zur Verfügung und wird dort auch von einer großen Entwicklergemeinschaft weiterentwickelt\cite[][9]{ng-book}. Da er unter der MIT-Lizenz veröffentlicht ist eignet sich AngularJS auch für den kommerziellen Einsatz\cites[9]{ng-book}[]{mit-license}.

Ende 2016 wurde eine neue Version von Angular veröffentlicht: Angular2\cite{ng2-out}. Durch die Bezeichnung kann AngularJS (Version 1) klar von Angular2 (Version 2) abgegrenzt werden. Im Rahmen dieser Bachelorarbeit wird AngularJS betrachtet.

\subsubsection{Funktionsweise}
Bei der Entwicklung von Webanwendungen mit AngularJS kommt das Model-View-ViewModel-Entwurfsmuster (MVVM), eine Erweiterung von Model-View-Controller (MVC), zum Einsatz. 

\subsubsection{Architektur}


\subsubsection{Dependency Injection}

\subsection{Node.js}
\subsubsection{Laufzeitumgebung}
Node.js ist eine Laufzeitumgebung, mit der JavaScript ohne Webbrowser ausgeführt werden kann\cite[][1]{hughes2012einfuehrungnodejs}. Somit ist es möglich, JavaScript nicht nur für die Darstellung von Benutzeroberflächen im Webbrowser zur Nutzen, sondern auch als Backend-Sprache oder zur Unterstützung von Entwicklungsprozessen auf Continuous-Integration-Servern oder Entwicklerarbeitsplätzen.

Intern nutzt Node.js die JavaScript-Engine \textit{Chrome V8} \cite{nodejs}, welche von Google als Open-Source-Software veröffentlicht wurde. V8 kommt auch im weitverbreiteten Webbrowser \textit{Google Chrome} zum Einsatz und implementiert den JavaScript-Standard ECMAScript wie in ECMA-262 spezifiziert\cite{chromev8}. ECMA-262 ist der im Juni 2016 veröffentlichte und zurzeit aktuellste JavaScript-Standard\cite{ecma262}. Somit bietet Node.js alle spezifizierten und von Google Chrome unterstützten Sprachfunktionalitäten. Es eignet sich daher auch für den Test von für Webbrowser entwickelte Webanwendungen.\todo{Unterschiede Node <-> JS}

\subsubsection{npm}
Der Node Package Manager (npm) ist der zusammen mit Node.js verteilte Paketmanager für JavaScript. Unter npm wird außerdem die \textit{npm Registry}, also die zentrale Ablage für JavaScript-Pakete, verstanden, auf welche der Node Package Manager zugreift. \cite{npm-about} Die npm Registry enthält über 180.000 Pakete und ist damit das größte Softwarerepository\cite{modulecount}.

Grundlegend funktioniert die Paketverwaltung mit einer JSON-Konfigurationsdatei, der \texttt{package.json} (vgl. Listing \ref{lst:package-json}). Die Datei enthält den Namen sowie die Version des Pakets, für welches sie angelegt wurde, sowie optional weitere Metadaten wie Beschreibung, Autor und Referenzlinks auf Bugtracker. Außerdem werden hier Abhängigkeiten angegeben, die zur Ausführung (\texttt{dependencies}) oder zur Entwicklung (\texttt{devDependencies}) in diesem Paket benötigt werden. \cite{npm-packagejson}Die angegebenen Abhängigkeiten werden von npm automatisiert heruntergeladen und im Ordner \texttt{node\_modules} abgelegt, von wo aus sie in die JavaScript-Anwendung eingebunden werden können. Auch transitive Abhängigkeiten werden von npm aufgelöst. \todo{Quelle}

\begin{figure}[H]
	\lstinputlisting[caption=Beispiel einer package.json, label=lst:package-json]{lst/package.json}
\end{figure}

\todo{Wo kommt die Quelle hin?}
Neben der Paketverwaltung kann npm auch zum Build als Taskrunner eingesetzt werden. Hiermit kann der Buildprozess eines Paketes automatisiert werden, z.\,B. durch die automatisierte Ausführung von Tests oder dem Aufrufen von Compilern. Hierzu werden in der package.json Skripte angegeben. Diese bestehen aus einem Skript-Namen und dem auszuführenden Befehl. Im Beispiel (siehe Listing \todo{}) werden drei Skripte definiert:\cite{cirkel-npmAsABuildTool}
\begin{itemize}
	\item \glqq lint\grqq~ führt das Kommando \texttt{jshint **.js} aus. Dies dient dem Überprüfen von JavaScript-Dateien auf statische Programmierfehler\cite{jshint-about}.
	\item \glqq build\grqq~ führt das Kommando \texttt{browserify [...]} aus. Dieses dient dem Zusammenfügen von mittels \texttt{require} eingebundenen JavaScript-Dateien in eine konkatenierte Datei\cite{browserify-about}.
	\item \glqq test\grqq~ führt das Kommando \texttt{mocha [...]} aus. Mocha ist ein Test-Runner (siehe auch \todo{Referenz, Quelle}).
\end{itemize}

Angegebene Skripte können auch automatisch in sogenannten Hooks (\textit{Pre} und \textit{Post Hooks}) ausgeführt werden. Im Beispiel (siehe Listing \ref{lst:packageWithScripts-json}) sind folgende Hooks definiert:\cite{cirkel-npmAsABuildTool}
\begin{itemize}
	\item \glqq prepublish\grqq~ wird vor der Ausführung von \texttt{publish}, welches ein npm Standard-Skript ist und das Paket in der npm Registry veröffentlicht\cite{npm-publish}, das benutzerdefinierte Skript \texttt{build} sowie \texttt{prebuild} und \texttt{postbuild} ausführen.
	\item \glqq prebuild\grqq~ wird vor Ausführung des \texttt{build}-Skripts das Paket durch Ausführung von \texttt{test} überprüfen.
	\item \glqq pretest\grqq~ wird vor Ausführung von \texttt{test} mittels \texttt{lint} das Paket auf statische Fehler untersuchen.
\end{itemize}

\begin{figure}[H]
	\lstinputlisting[caption={Beispiel von Skripten in einer package.json (aus \cite{cirkel-npmAsABuildTool})}, label=lst:packageWithScripts-json]{lst/packageWithScripts.json}
\end{figure}

Die wohl populärste Funktion von Taskrunnern ist das automatisierte Beobachten des Dateisystems auf Änderungen. Häufig ist es wünschenswert, dass bei einer Dateiänderung automatisch ein entsprechender Buildprozess oder die Tests ausgeführt werden. Diese Funktionalität bietet npm im Gegensatz zu anderen Taskrunnern wie \textit{Gulp} oder \textit{Grunt} nicht nativ, sondern nur mithilfe eines Pakets. Ein entsprechendes Beispiel, welches bei Veränderung einer Datei im Paket-Ordner die JavaScript-Module zu einer Datei konkateniert\cite{browserify-about} und den JavaScript-Code des Paketes überprüft, findet sich in Listing \ref{lst:packageWithWatch-json}.\cite{cirkel-npmAsABuildTool}

\begin{figure}[H]
	\lstinputlisting[caption={Beispiel der Watch-Funktionalität in einer package.json (aus \cite{cirkel-npmAsABuildTool}, angepasst durch den Autor)}, label=lst:packageWithWatch-json]{lst/packageWithWatch.json}
\end{figure}

%\newpage
%\section{Systeminformationen}\label{system_info}

%\newpage
%\section{Rahmenbedingungen}\label{rahmenbedinugnen}

%\newpage
%\section{Durchführung}\label{durchfuehrung}

%\newpage
%\section{Ergebnisse}\label{ergebnisse}

%\newpage
%\section{Abschluss}\label{abschluss}

%\onecolumn
% einfacher Zeilenabstand
%\singlespacing
% Literaturliste soll im Inhaltsverzeichnis auftauchen
%\newpage
%\addcontentsline{toc}{section}{Literaturverzeichnis}
% Literaturverzeichnis anzeigen
%\renewcommand\refname{Literaturverzeichnis}
%\bibliography{BachelorArbeit}

%% Index soll Stichwortverzeichnis heissen
% \newpage
% % Stichwortverzeichnis soll im Inhaltsverzeichnis auftauchen
% \addcontentsline{toc}{section}{Stichwortverzeichnis}
% \renewcommand{\indexname}{Stichwortverzeichnis}
% % Stichwortverzeichnis endgueltig anzeigen
% \printindex

\onehalfspacing
% evtl. Anhang
%\newpage
%\addcontentsline{toc}{section}{Anhang}
%\fancyhead[L]{Anhang} %Kopfzeile links
%%\{Anhang}\label{anhang}

% Eidesstattliche Erklärung
%\addcontentsline{toc}{section}{Eidesstattliche Erklärung}
%\section*{Eidesstattliche Erklärung}
\thispagestyle{empty}

\begin{verbatim}

\end{verbatim}

\begin{LARGE}Eidesstattliche Erklärung zur Bachelorarbeit\end{LARGE}
\begin{verbatim}


\end{verbatim}
Ich versichere, die von mir vorgelegte Arbeit selbstständig verfasst zu haben. Alle Stellen, die wörtlich oder sinngemäß aus veröffentlichten oder nicht veröffentlichten Arbeiten anderer entnommen sind, habe ich als entnommen kenntlich gemacht. Sämtliche Quellen und Hilfsmittel, die ich für die Arbeit benutzt habe, sind angegeben. Die Arbeit hat mit gleichem Inhalt bzw. in wesentlichen Teilen noch keiner anderen Prüfungsbehörde vorgelegen.



\begin{displaymath}
% use packages: array
\begin{array}{ll}
Unterschrift:~~~~~~~~~~~~~~~~~~~~~~~~~~~~~~~~~~~~~~~~~~
& Ort, Datum:~~~~~~~~~~~~~~~~~~~~~~~~~~~~~~~~~~~~~~~~~~
\end{array}
\end{displaymath}


% leere Abschlussseite
%\newpage
%\thispagestyle{empty} % erzeugt Seite ohne Kopf- / Fusszeile
\section*{ }

\end{document}